\subsection{Initial-Condition Route for Every Solver Task}

An initial condition (IC) specifies the matter distribution in position and
velocity, its total mass and units, and the starting field or metric. A grid
solver represents this distribution by values or averages on a phase space
grid. A marker solver represents it by weighted numerical particles, called
markers. A restart also retains the evolved time, event clocks,
random-number counters, and method-specific internal state. Those additional
quantities let the solver continue the same numerical trajectory.

The rows are ordered from the maintained production marker methods to the
grid and research methods.  Each row names both the public solver executable
and the source file that actually reads or constructs its starting state.
An executable ending in \path{_openmp} uses the same IC route as the named CPU
task; \path{_openmp} means that loop work may be shared among CPU threads.
The two grid executables ending in \path{_mpi} likewise use the same route;
\path{_mpi} means that the grid is split among communicating processes.

\begin{longtable}{@{}p{0.28\linewidth}p{0.28\linewidth}p{0.36\linewidth}@{}}
\toprule
Consuming task & Fresh-data or restart route & Authoritative implementation
and required checks \\
\midrule
\endfirsthead
\toprule
Consuming task & Fresh-data or restart route & Authoritative implementation
and required checks \\
\midrule
\endhead
\path{venus_cpu_vlasov_poisson_wpic_expanding} & Analytic expanding power-law
markers generated inside the driver, or a versioned weighted particle-in-cell
(WPIC) restart selected by \code{--restart} & The task main is
\path{cpu/src/apps/vlasov_poisson/expanding_wpic.cpp};
\path{cpu/src/drivers/expanding_wpic.cpp} generates, reads, and validates the
markers.  Check marker--mesh coverage, expansion and velocity conventions,
represented mass, and restart clocks. \\

\path{venus_cpu_subhalo_wpic} & Analytic expanding power-law markers or a
complete WPIC restart carrying masses, internal state, IDs, and event clocks &
The task main is \path{cpu/src/apps/subhalo/expanding_wpic.cpp}; IC generation
and restart loading are in \path{cpu/src/drivers/expanding_wpic.cpp}, and the
collision state is defined in
\path{cpu/include/ev/weighted_vlasov_pic_subhalo_collisions.hpp}.  Check the
physical subhalo masses, Plummer calibration, locality, pair majorant,
capture ledger, and collisionless control. \\

\path{venus_cpu_sidm_wpic_constant} & Exactly one of
\code{--frank-restart PATH}, which imports a persistent Frank SphDSMC particle
population, or \code{--pic-restart PATH}, which resumes a WPIC state; there is
no separate analytic-fresh option & The task main is
\path{cpu/src/apps/sidm/wpic_constant.cpp}.  The command-line check and import
are in \path{cpu/src/drivers/sidm_wpic.cpp}; the conversion itself is in
\path{cpu/include/ev/weighted_pic_frank_restart_import.hpp}.  Check positive
marker masses, the static-halo convention, shell locality, collision clocks,
and any equal-mass assumption. \\

\path{venus_cpu_sidm_wpic_velocity_dependent} & The same exactly-one-of
\code{--frank-restart} or \code{--pic-restart} route as the constant WPIC task
& The task main is \path{cpu/src/apps/sidm/wpic_velocity_dependent.cpp}; the
reader and checks are again in \path{cpu/src/drivers/sidm_wpic.cpp}.
Check the selected power-law or Yukawa--Born formula over the realized
relative-speed range; the task is a research extension rather than a
paper-exact constant-law claim. \\

\path{venus_cpu_sidm_sphdsmc_constant} and
\path{venus_cpu_sidm_sphdsmc_constant_openmp} & Sample an analytic or
\code{EVICRWL1} phase-grid distribution into persistent markers, or resume a
versioned spherical direct-simulation Monte Carlo (SphDSMC) restart & The task
main is \path{cpu/src/apps/sidm/sphdsmc_constant.cpp}; loading, sampling, and
restart checks are in \path{cpu/include/ev/spherical_dsmc_solver.hpp}.
\path{cpu/tools/generate_frank_core_spherical_dsmc_restart.py} and
\path{cpu/tools/generate_frank_nfwc40_fig2_restart.py} construct the two named
benchmark populations.  Check marker masses and identifiers, radial support,
equilibrium, shell occupancy, and mean-free-path locality. \\

\path{venus_cpu_sidm_sphdsmc_velocity_dependent} and
\path{venus_cpu_sidm_sphdsmc_velocity_dependent_openmp} & The same
persistent-marker routes as the constant task & The task main is
\path{cpu/src/apps/sidm/sphdsmc_velocity_dependent.cpp}; the common reader is
\path{cpu/include/ev/spherical_dsmc_solver.hpp}.  Recheck the cross-section
formula, sampled velocity range, collision-rate upper bound, locality, and
long-time behavior; validation of the constant law does not validate this
extension. \\

\path{venus_cpu_einstein_vlasov_wpic_low_compactness} & An internally
generated paired Einstein--Vlasov/Vlasov--Poisson marker realization, or
\code{--resume} from the two-slot restart named by
\path{low_restart_manifest.txt} in the same output directory & The task main is
\path{cpu/src/apps/einstein_vlasov/wpic_low_compactness.cpp}; construction and
restart selection are in \path{cpu/src/drivers/einstein_vlasov_wpic.cpp}.
Check fixed-point marker weights, identical marker IDs in the pair, low
compactness, and the normalization used in the comparison. \\

\path{venus_cpu_einstein_vlasov_wpic_black_hole} with
\code{--coordinates schwarzschild} & Internally generated Plummer, shell, or
bound energy-cutoff markers; a coordinate-neutral physical-IC artifact; or
\code{--resume} from \path{high_einstein_vlasov.restart} in the output
directory & The task main is
\path{cpu/src/apps/einstein_vlasov/wpic_black_hole.cpp}; the driver is
\path{cpu/src/drivers/einstein_vlasov_wpic.cpp}, and the energy-cutoff sampler
is \path{cpu/include/ev/weighted_einstein_vlasov_static_shell.hpp}.  Check the
continuum parameters, total mass rescaling, kick, angular scale, quiet
sampling, marker IDs, binding diagnostics, and exterior metric. \\

\path{venus_cpu_einstein_vlasov_wpic_black_hole} with
\code{--coordinates eddington_finkelstein} & A fresh run must specify
\code{--ef-null-initial-data PATH}; a raw constant-Schwarzschild-time marker
file is not enough.  Continuation uses \code{--resume}, the EF checkpoint in
the output directory, and its matching post-horizon audit sidecar & The shared
driver, \path{cpu/include/ev/weighted_einstein_vlasov_coordinate_map.hpp}, and
\path{cpu/tools/validate_wpic_null_map.py} own the import and checks.  Require
donor format version 2, exactly one crossing per marker ID, face-aware field
reconstruction, clearance between matter and the map boundary, convergence of
the coordinate map, all nested file hashes, and fixed null-slice mass
\(M_{\rm null}\). \\

\path{venus_cpu_einstein_vlasov_eulerian}, its \path{_openmp} executable, and
\path{venus_cpu_einstein_vlasov_eulerian_mpi} & Analytic distribution or
\code{EVICRWL1} phase-grid file; versioned grid restart for continuation & The
task main is \path{cpu/src/apps/einstein_vlasov/eulerian.cpp};
\path{cpu/include/ev/initial_data.hpp} reads grid data, while
\path{cpu/python/ev_solver/ic_generator.py} and
\path{cpu/tools/ic_generator_legacy.py} generate general files.  Check that
the distribution is nonnegative, the file grid matches the run grid, mass is
normalized, matter stays clear of unsupported boundaries, and the initial
metric constraints close. \\

\path{venus_cpu_vlasov_poisson_eulerian}, its \path{_openmp} executable, and
\path{venus_cpu_vlasov_poisson_eulerian_mpi} & Analytic or \code{EVICRWL1}
phase-grid data; versioned grid restart & The task main is
\path{cpu/src/apps/vlasov_poisson/eulerian.cpp}; the reader and generators are
the same as for grid Einstein--Vlasov.  An expanding run must also record the
scale factor, time origin, Poisson source, and whether stored velocities are
physical or measured relative to the expansion. \\

\path{venus_cpu_boltzmann_poisson_finite_angle},
\path{venus_cpu_boltzmann_poisson_exact},
\path{venus_cpu_boltzmann_poisson_reversible},
\path{venus_cpu_boltzmann_poisson_reversible_sampled}, and
\path{venus_cpu_boltzmann_poisson_reversible_streamed}, including each
\path{_openmp} executable & Eulerian phase-grid data or versioned grid restart
& Their task mains are the correspondingly named files under
\path{cpu/src/apps/boltzmann_poisson};
\path{cpu/include/ev/initial_data.hpp} supplies \(f(r,w,\eta)\), the stored
distribution in radius, radial velocity, and half the squared specific
angular momentum, \(\eta=L^2/2=r^2v_t^2/2\), where \(v_t\) is tangential
speed.  Reversible tasks also require a compatible cell-average grid,
closure under reversal of radial velocity, full collision support, and valid
projection and omission limits. \\

\path{venus_cpu_boltzmann_poisson_expanding_prompt_cusp} and its
\path{_openmp} executable & The supplied configuration reads the bundled
\code{EVICRWL1} zero-flux file
\path{cpu/examples/initial_data/bp_zero_flux_gamma2p5_eds_R20_rmin0p1_W3_L8_Nr48_Nw32_NL24_quadL_wstretch3p5_discrete.bin};
a compatible grid file or restart may replace it & The task main is
\path{cpu/src/apps/boltzmann_poisson/expanding_prompt_cusp.cpp}; the INI interface selects the grid distribution and collision controls.
There is no dedicated prompt-cusp IC generator in this repository.  Check the
finite radial box, expansion convention, and density-contrast source. \\

\path{venus_cpu_boltzmann_poisson_fokker_planck} and
\path{venus_cpu_boltzmann_poisson_landau}, including their \path{_openmp}
executables & Eulerian Vlasov--Poisson/Boltzmann--Poisson phase-grid data or a
grid restart & Their task mains are the correspondingly named files under
\path{cpu/src/apps/boltzmann_poisson}; \path{cpu/include/ev/initial_data.hpp}
reads the state, and
\path{cpu/tools/generate_frank_core_cell_average_restart.py} builds the Frank
grid benchmark.  Check the background distribution used by the approximation,
velocity support, collision-law units, and numerical coefficient floors. \\

\path{venus_cpu_boltzmann_poisson_canonical_orbit} and its \path{_openmp}
executable & \code{--restart PATH} is always required and supplies a versioned
Eulerian grid restart.  The task also always loads a stationary Eddington table:
\code{--eddington-table PATH}, or by default
\path{frank_core_eddington_8192.csv} beside that source restart.
\code{--canonical-restart PATH} optionally restores the evolved orbit atlas
and event state; it does not replace either required source input & The task
main and parser are
\path{cpu/src/apps/boltzmann_poisson/canonical_orbit.cpp};
\path{cpu/include/ev/canonical_boltzmann_poisson_restart.hpp} owns the
canonical state, and \path{cpu/tools/export_frank_eddington_table.py} writes
the table.  Validate atlas coverage, packet mass, random counters, event
clocks, and restart equivalence. \\

Python Einstein--Vlasov reference package & Built-in analytic shell,
\code{EVICRWL1} grid data, or a compatible \code{EVLOGRST1} restart &
\path{cpu/python/ev_solver/solver_setup.py} selects the route,
\path{cpu/python/ev_solver/initial_data.py} builds or reads fresh data, and
\path{cpu/python/ev_solver/io.py} validates the restart.  This package is an
EV reference only: it rejects VP, adaptive-\(w\), non-KDK characteristic
integrators, and the C++-only analytic options named by its validator. \\

\path{venus_gpu_vlasov_poisson} and
\path{venus_gpu_vlasov_poisson_mpi} & Mandatory versioned CPU/GPU
WPIC restart; no GPU-native fresh-data generator &
The task main is \path{gpu/tools/venus_gpu_vlasov_poisson.cu};
\path{gpu/src/wpic_production_driver.cu} runs the task, and
\path{gpu/src/wpic_restart.cpp} reads and uploads the state.  Validate restart
schema, marker identity, mesh contract, CPU/GPU short-run parity, and the
requested rank layout. \\

\path{venus_gpu_subhalo_boltzmann_poisson} and
\path{venus_gpu_subhalo_boltzmann_poisson_mpi} & Mandatory WPIC
restart carrying the subhalo state &
The task main is \path{gpu/tools/venus_gpu_subhalo_boltzmann_poisson.cu}; the
common driver, \path{gpu/src/wpic_subhalo_collisions.cu}, and
\path{gpu/src/wpic_restart.cpp} load it.  Check CPU calibration provenance,
internal state, cell ownership and migration, pair invariants, and restart
equivalence. \\

\path{venus_gpu_sidm_boltzmann_poisson},
\path{venus_gpu_sidm_constant}, \path{venus_gpu_sidm_power_law}, and
\path{venus_gpu_sidm_yukawa_born}, together with the four full names ending in
\path{_mpi} & Mandatory static-halo WPIC restart; there is no GPU-native fresh
generator & Their four task mains are the correspondingly named files under
\path{gpu/tools}; \path{gpu/src/wpic_production_driver.cu},
\path{gpu/src/wpic_collisions.cu}, and \path{gpu/src/wpic_restart.cpp} load the
state.  Check the selected law, units, relative-speed support, local shell
clock, locality, and CPU/GPU convergence.  The current material-collision MPI
path is not invariant under a change in rank count. \\
\bottomrule
\end{longtable}

Check a generated state against the consuming solver's physical and file
conventions. Preserve the generator command, input distribution or continuum
parameters, random or low-discrepancy realization, normalization report,
output hash, and import checks with the run. These records establish which
distribution the solver actually received.


\paragraph{A bound Gaussian shell for GPU Vlasov--Poisson.}
The CPU utility \path{cpu/tools/generate_bound_gaussian_shell_restart.py}
constructs the two background representations of the same physical shell.
Its density is proportional to \(\exp[-(r-10)^2/2]\) on \(7<r<13\).
A radial interval contains its density multiplied by the spherical shell
volume. The cumulative mass distribution used to draw radii therefore contains
the squared-radius factor:
\[
 F_r(r)=\frac{\int_7^r s^2\exp[-(s-10)^2/2]\,\dd s}
                  {\int_7^{13} s^2\exp[-(s-10)^2/2]\,\dd s},
 \qquad 7<r<13,
\]
where \(F_r\) is the fraction of the shell mass inside radius \(r\).
The generator inverts this cumulative distribution at sample points chosen
to cover the unit interval evenly. Each sampled radius receives two markers
with opposite physical radial velocities, so the pair has zero mean radial
motion. Their Gaussian scale is \(\sigma_v=0.08481562072248405\),
with \(|u|<3\sigma_v\). Two Gaussian transverse velocity components give a
Rayleigh distribution of transverse speed \(v_t\), truncated at
\(v_t<3\sigma_v\). The stored angular variable is \(\eta=r^2v_t^2/2\).
The default \(1048576\) equal-mass markers represent total mass
\(3.5968447592701334\). The generator reconstructs their spherical potential
and verifies that every initial specific energy
\((u^2+v_t^2)/2+\Phi\) is negative.

The static file stores \(w=u\). The expanding file stores \(w=u-Hr\), with
initial \(H=1/18\). Both files thus describe identical physical positions,
velocities, masses, angular momenta and particle identifiers. Their fresh
version-6 particle states carry zero event ages and infinite initial timescale
estimates. The solver constructs the initial gravitational field on import.
\codesource{Bound Gaussian-shell initial data}{\path{cpu/tools/generate_bound_gaussian_shell_restart.py}}{Samples the radial mass distribution and Gaussian velocities, checks binding, and writes static and expanding particle states with their mesh sidecars}
\begin{lstlisting}[language=bash]
python3 cpu/tools/generate_bound_gaussian_shell_restart.py \
  --output-directory output/gaussian_ic --particles 1048576
\end{lstlisting}
Run this command from the repository root with a new output directory.
The two solver commands in
\path{cpu/examples/modular/vp_bound_gaussian_shell_gpu.md} select the
corresponding restart, background and 512-cell logarithmic mesh.
They use single-GPU Vlasov--Poisson with a fixed KDK cap of \(0.0075\)
and an outer reflecting radius of \(10000\).

\paragraph{A user-defined grid distribution.}
The supplied \path{cpu/examples/user_distribution.py} defines the vectorized
function \code{distribution(r,w,L)} for a compact inward-moving shell.
The companion \path{cpu/examples/modular/vp_custom_shell.ini} selects a
static Vlasov--Poisson calculation. The repository README gives the exact
commands to evaluate the function into an \code{EVICRWL1} file, run the
solver, and plot the resulting tables. The binary supplies its coordinate
arrays; the INI file supplies the equation, boundaries, integration controls,
and output cadence. Paths in that INI file are relative to its own directory.
The example writes all generated input and output under
\path{output/cpu/custom_shell}. Its distribution function can be edited
without changing a solver source file.
